% 对应题号 M1-14；图与本题条件、推导、结论逐项对应。
% 图契约：印刷宽16cm；最小字号8.4pt；几何/场或曲线在左，计算关系在右。
\documentclass[tikz,border=7pt]{standalone}
\usepackage{fontspec,xeCJK,amsmath,amssymb,graphicx}
\IfFontExistsTF{Times New Roman}{\setmainfont{Times New Roman}}{\setmainfont{TeX Gyre Termes}}
\IfFontExistsTF{Heiti SC}{\setCJKmainfont[AutoFakeBold=1.5]{Heiti SC}}{\setCJKmainfont[AutoFakeBold=1.5]{FandolHei-Regular}}
\usetikzlibrary{arrows.meta,bending,calc,patterns,positioning}
\definecolor{ink}{HTML}{243745}
\definecolor{blue}{HTML}{275DAB}
\definecolor{green}{HTML}{177C63}
\definecolor{red}{HTML}{BD4935}
\definecolor{orange}{HTML}{AA6C16}
\definecolor{purple}{HTML}{715393}
\tikzset{every node/.style={font=\fontsize{9.3}{12.5}\selectfont,text=ink},
 axis/.style={-{Stealth[length=4pt]},line width=.65pt,draw=ink!65},
 vec/.style={-{Stealth[length=5pt]},line width=1pt,draw=blue},
 short/.style={-{Stealth[length=3pt]},line width=.8pt},
 guide/.style={dashed,draw=ink!40,line width=.5pt},
 card/.style={rounded corners=2pt,draw=blue!65,fill=blue!4,align=center,text width=6.6cm,inner xsep=9pt,inner ysep=7pt},
 title/.style={font=\fontsize{10.5}{14}\selectfont\bfseries,text width=15cm,align=center},
 label/.style={font=\fontsize{8.4}{11}\selectfont,align=center}}
\begin{document}
\begin{tikzpicture}[x=1cm,y=1cm]
\path[use as bounding box] (-8,-3.4) rectangle(8,3.4);
\node[title] at (0,3) {M1-14\quad 本题的等值线、梯度与旋度方向};
\begin{scope}[shift={(-4,0)}]

\draw[axis](-2.05,.2)--(1.7,.2)node[right]{$x$};
\draw[axis](-1.8,0)--(-1.8,2.05)node[above]{$y$};
\foreach\c in{.4,1,1.6}{\pgfmathsetmacro{\lo}{sqrt(\c/1.5)}
\draw[blue!65]plot[domain=\lo:2.2,samples=70]({-1.8+1.2*\x},{.2+1.2*\c/(\x*\x)});}
\coordinate(P)at(-.6,1.4);\fill[red](P)circle(2pt)node[below left=2pt]{$P$};
\draw[vec,green](P)--(.48,1.94);
\node[label,green,fill=white,inner sep=1.5pt]at(.95,2.15){$\nabla f(P)=(2,1,0)$};
\node[label]at(-.15,2.62){$z=0$：$f=x^2y$；$P=(1,1,0)$};
\node[label]at(1.65,1.1){等值线\\$x^2y=C$};

\draw[axis](-1.7,-1.7)--(1.3,-1.7)node[right]{$x$};
\draw[axis](-1.5,-1.9)--(-1.5,-.3)node[above]{$y$};
\coordinate(Q)at(-.7,-.9);\fill[ink](Q)circle(2pt)node[below=4pt]{$P$};
\draw[vec,green](Q)--(.1,-.9)node[right]{$\mathbf A(P)$};
\draw[vec,red](Q)--(-1.5,-.9);
\node[label,red,fill=white,inner sep=1pt]at(-.75,-.47){$(\nabla\times\mathbf A)(P)=-\mathbf e_x$};
\draw[vec,blue!60](-1.1,-1.38)--(-.97,-1.38);
\draw[vec,blue!60](.05,-1.38)--(.62,-1.38);
\node[label]at(-.1,-2.18){$z=0:\ \mathbf A=x^2\mathbf e_x$};
\node[label]at(-.1,-2.65){$\mathbf A(P)=\mathbf e_x$，旋度在本切片内};

\end{scope}\node[card,draw=blue!65,fill=blue!4] at (3.95,1.65) {$f=x^2y+z^2$\\$\mathbf A=(x^2,yz,xz)$};
\node[card,draw=green!65,fill=green!4] at (3.95,0) {$\nabla f=(2xy,x^2,2z)$\\$\nabla^2f=2y+2$};
\node[card,draw=orange!65,fill=orange!4] at (3.95,-1.65) {$\nabla\cdot\mathbf A=3x+z$\\$\nabla\times\mathbf A=(-y,-z,0)$};
\end{tikzpicture}\end{document}
